\documentclass{article}
\usepackage{amsmath, amssymb}

\begin{document}

\section*{1. Tails -- Definition \& Motivation}

\subsection*{Definition}
The tail of a random variable $X$ is:
\[
P(X > x)
\]
This represents the probability of extreme values (right side of distribution).

\subsection*{Why do we care about tails?}
Examples:
\begin{itemize}
    \item Jobs delayed more than 24 hours
    \item Packets dropped due to full buffer
\end{itemize}

These are rare but critical events.

\subsection*{Key Idea}
\begin{itemize}
    \item Mean $\rightarrow$ average behavior (easy)
    \item Tail $\rightarrow$ rare extreme events (hard)
\end{itemize}

\subsection*{Why tails are hard?}
Requires summing many small probabilities.

\section*{2. Tail Example (Server Problem)}

\subsection*{Problem}
Suppose $n$ jobs are distributed among $n$ servers randomly. Find:
\[
P(\text{a server gets } \ge k \text{ jobs})
\]

\subsection*{Intuition}
\begin{itemize}
    \item Jobs assigned randomly
    \item Most servers get few jobs
    \item Occasionally one server gets many jobs
\end{itemize}

This is a tail event (rare overload).

\subsection*{Model}
\[
X \sim \text{Binomial}(n, p)
\]

\subsection*{Exact Probability}
\[
P(X \ge k) = \sum_{i=k}^{n} \binom{n}{i} p^i (1-p)^{n-i}
\]

This is hard to compute, motivating inequalities.

\section*{3. Markov's Inequality}

\subsection*{Definition}
For a non-negative random variable $X$:
\[
P(X \ge a) \le \frac{E[X]}{a}
\]

\subsection*{Assumptions}
\begin{itemize}
    \item $X \ge 0$
    \item $a > 0$
\end{itemize}

\subsection*{Key Idea}
Bounding probability using expectation.

\subsection*{Interpretation}
\begin{itemize}
    \item Gives an upper bound on tail probability
    \item Works with minimal information (only mean)
\end{itemize}

\section*{4. Chebyshev's Inequality}

\subsection*{Definition}
\[
P(|X - \mu| \ge k\sigma) \le \frac{1}{k^2}
\]

\subsection*{Assumptions}
\begin{itemize}
    \item $\mu = E[X]$
    \item $\sigma^2 = \text{Var}(X)$
\end{itemize}

\subsection*{Key Idea}
Uses variance to bound deviation from mean.

\subsection*{Interpretation}
\begin{itemize}
    \item Gives probability of large deviations
    \item Stronger than Markov
\end{itemize}

\section*{5. Chernoff Bound}

\subsection*{Core Idea}
\[
P(X \ge a) = P(e^{tX} \ge e^{ta})
\]

\subsection*{Lower Tail}
For $t < 0$:
\[
P(X \le a) = P(e^{tX} \ge e^{ta})
\]

\[
P(X \le a) \le \min_{t < 0} \left\{ \frac{E[e^{tX}]}{e^{ta}} \right\}
\]

\subsection*{Poisson Tail Derivation}

Let $X \sim \text{Poisson}(\lambda)$.

\subsubsection*{Step 1: Compute $E[e^{tX}]$}
\[
E[e^{tX}] = \sum_{i=0}^{\infty} e^{ti} \cdot \frac{e^{-\lambda} \lambda^i}{i!}
\]

\[
= e^{-\lambda} \sum_{i=0}^{\infty} \frac{(\lambda e^t)^i}{i!}
\]

\[
= e^{-\lambda} \cdot e^{\lambda e^t}
\]

\[
= e^{\lambda(e^t - 1)}
\]

\subsubsection*{Step 2: Bound Tail}
For $a > \lambda$:
\[
P(X \ge a) \le \min_{t > 0} \left\{ e^{\lambda(e^t - 1) - ta} \right\}
\]

\subsubsection*{Step 3: Optimize $t$}
\[
t = \ln\left(\frac{a}{\lambda}\right)
\]

\subsubsection*{Final Result}
\[
P(X \ge a) \le e^{\lambda\left(\frac{a}{\lambda} - 1\right) - a \ln\left(\frac{a}{\lambda}\right)}
\]

\subsubsection*{Simplified Form}
\[
P(X \ge a) \le e^{a - \lambda} \left(\frac{\lambda}{a}\right)^a
\]

\section*{6. Jensen's Inequality}

\subsection*{Convex Function}
A function $f$ is convex if:
\[
f(E[X]) \le E[f(X)]
\]

\subsection*{Jensen's Inequality}
For convex $f$:
\[
f(E[X]) \le E[f(X)]
\]

\subsection*{Key Idea}
Expectation interacts with convexity.

\section*{7. Case Study}

\subsection*{Problem}
How to predict worst-case demand on a server?

\subsection*{Insight}
\begin{itemize}
    \item Use tail bounds
    \item Exact computation is hard
    \item Inequalities give safe upper bounds
\end{itemize}

\end{document}